\subsection{Separation into blocks}

Evaluate at $\boldsymbol{q}=1$ to obtain the ordinary multiplicities
\[
t_{j,p}(1) = \sum_{d=0}^{N_p} f_{p,d}^{j}.
\]

The simple modules are partitioned into blocks according to the non-vanishing of these ordinary multiplicities. Starting with the smallest unclassified index $j$, we define

\[
B_j = \{p : t_{j,p}(1) \neq 0\}.
\]

The procedure is continued until every simple module belongs to exactly one block.

\begin{table}[ht]
\centering
\begin{tabular}{c|c|l}
\text{Block} & \text{Indices} & \text{Simple }D(G)\text{-modules} \\\hline
$B_{1}$ & $\{1,\ 5\}$ & $M(e,+)$,\ $M(\sigma,-)$ \\
$B_{2}$ & $\{2\}$ & $M(e,-)$ \\
$B_{3}$ & $\{3,\ 6\}$ & $M(e,\rho)$,\ $M(\tau,0)$ \\
$B_{4}$ & $\{4\}$ & $M(\sigma,+)$ \\
$B_{5}$ & $\{7\}$ & $M(\tau,1)$ \\
$B_{6}$ & $\{8\}$ & $M(\tau,2)$ \\
\end{tabular}
\caption{Block decomposition of the simple $D(G)$-modules.}
\label{tab:block-decomposition}
\end{table}

